\section{Introdução: A Distância entre a Promessa e a Realidade}

A promessa transformacional dos gêmeos digitais na odontologia, minuciosamente delineada nos capítulos anteriores, contrasta de forma contundente com um cenário de adoção ainda incipiente e permeado por incertezas estruturais. Se, por um lado, os fundamentos tecnológicos da modelagem 3D e da simulação in silico (\autoref{ch:fundamentos}), as instigantes aplicações clínicas (\autoref{ch:ortodontia}, \autoref{ch:implantodontia}, \autoref{ch:endodontia}) e o potencial de ecossistemas abertos (\autoref{ch:opencode}) demonstram um horizonte revolucionário, por outro, a distância ontológica e operacional entre a prova de conceito (Proof of Concept - PoC) e a prática clínica generalizada permanece considerável.

Este capítulo adota uma postura deliberadamente crítica, distanciando-se do determinismo tecnológico que frequentemente impregna o discurso sobre inovações na saúde. Em vez de celebrar antecipadamente os avanços, o texto examina as fraturas estruturais que impedem a maturação da tecnologia e a sua translação efetiva para a beira do equipo. A análise organiza-se em seis dimensões fundamentais — técnica, clínica, regulatória, econômica, de pesquisa e comparativa — com o escopo de transformar os obstáculos identificados em um roteiro estratégico para a ação coordenada do setor. 

Como advertem Duggal et al. \cite{duggal2026exploring}, a entusiasmada literatura sobre gêmeos digitais odontológicos, que cresceu exponencialmente nos últimos anos, carece tragicamente de validação empírica robusta. A vasta maioria das publicações restringe-se a discussões teóricas ou pequenos estudos in vitro, e é precisamente esta lacuna translacional que este capítulo busca escrutinar em profundidade, fornecendo alicerces para o desenvolvimento maduro e responsável da odontologia baseada em simulações.
